\section{Conjuntos bem-ordenáveis}\label{sec:bemOrdenaveis}
\index{função escolha|see{função de escolha}}\index{função de escolha}
Suponha que exista uma função escolha $f$ para $\mathcal P(A)\setminus{\emptyset}$.
Intuitivamente, podemos, via recursão transfinita, bem-ordenar $A$:
vamos, um por vez, etiquetando os elementos de $A$ com ordinais $0, 1, 2, \dots \omega, \omega+1, \dots, \omega\cdot 2, \dots$, até exaurir os elementos de $A$.
Automaticamente, tal recursão nos dá uma bijeção de $A$ com um ordinal, o que induz, em $A$, uma boa-ordem.

Um argumento informal para defender a validade do Axioma da Escolha é o seguinte: dado uma coleção de conjuntos não vazios, como montar uma função que seleciona exatamente um elemento de cada?
Ora, ``selecionando qualquer um''.
O problema com esse argumento é que no nosso formalismo, não há em nossa linguagem uma ferramenta que nos permite ``selecionar qualquer um'' e rotular tais elementos para se obter tal função.
Intuitivamente, podemos pensa em pegar ``o primeiro elemento que vier à mente'' de cada conjunto, mas, novamente, não há uma ferramenta formal que nos permita fazer isso.
Porém, se temos uma boa ordem na reunião dos conjuntos em que estamos tentando selecionar tais elementos, tal raciocínio é formalizável: fixe uma tal boa ordem e tome, sim, o ``primeiro elemento de cada conjunto'' com respeito a tal boa ordem.

A formalização destes raciocínio nos dá a seguinte proposição.

\begin{proposition}[$ZF^-$]\label{prop:funcaoEscolhaImplicaBoaOrdem}
    Seja $A$ um conjunto.
    São equivalentes:
    \begin{enumerate}[label=(\alph*)]
        \item Existe uma função escolha para $\mathcal P(A)\setminus\{\emptyset\}$.
        \item $A$ é bem-ordenável.
    \end{enumerate}
\end{proposition}
\begin{proof}
    (b) $\Rightarrow$ (a): caso $A$ seja bem-ordenável, seja $R$ uma boa-ordem em $A$.
    Defina $c:\mathcal P(A)\setminus\{\emptyset\}\to A$ dada por:
    \[
        c(X) = \min_R X
    \]
    É imediato que $c$ é uma função escolha para $\mathcal P(A)\setminus\{\emptyset\}$.

    Reciprocamente, suponha que exista uma função escolha $c$ para $\mathcal P(A)$.
    Fixe um elemento $b \notin A$.

    Recursivamente, introduza um símbolo funcional $F:\ON \to A$ tal que, para cada ordinal $\alpha$,
    \[
        F(\alpha)=\begin{cases}
        c\left(A\setminus \{F(\beta): \beta<\alpha\}\right), & \text{se } A\setminus \{F(\beta): \beta<\alpha\}\neq \emptyset\\
        b, & \text{caso contrário}
        \end{cases}
    \]

    Existe um ordinal $\alpha$ tal que $F(\alpha)=b$.
    Do contrário, para todo ordinal $\alpha$, $F(\alpha)\in A\setminus\{F(\beta): \beta<\alpha\}$, de modo que $F$ é injetora.
    Em particular, $F|h(A): h(A)\to A$ é injetora, em que $h(A)$ é o número de Hartogs de $A$, um absurdo.

    Seja $\delta$ o menor ordinal tal que $F(\delta)=b$.
    Temos que $F| \delta: \delta \to A$ é uma função injetora pelo mesmo argumento acima, e $F(\delta)=b\notin A$, de modo que $A=\{F(\beta): \beta<\delta\}$.
    Ou seja, $F|\delta$ é uma bijeção de $\delta$ em $A$.

    Dessa forma, $A$ é bem-ordenável, pois a bijeção $F|\delta$ de $\delta$ em $A$ induz uma boa-ordem em $A$.
\end{proof}


Assim, temos o seguinte corolário.
\begin{corollary}[$ZF^-$]\label{cor:equivalenciaFuncaoEscolhaBoaOrdem}
    São equivalentes.
    \begin{enumerate}[label=(\alph*)]
        \item O Axioma da Escolha.
        \item Todo conjunto é bem-ordenável.
    \end{enumerate}
\end{corollary}

A definição sobre boa ordenação, apesar de menos intuitiva, é mais simples: seu enunciado não depende do Axioma da Potência.
Por isso, neste texto, oficialmente, adotaremos ela como o Axioma da Escolha.

\begin{axiom}\label{ax:escolha}\index{axioma!da escolha}
    Todo conjunto é bem-ordenável.
\end{axiom}